\documentclass[11pt,a4paper,fontset=fandol]{ctexart}

\usepackage[a4paper,top=2.25cm,bottom=2.45cm,left=2.25cm,right=2.25cm,headheight=18pt,headsep=0.55cm,footskip=24pt]{geometry}
\usepackage{amsmath,amssymb,amsthm}
\usepackage{fancyhdr}
\usepackage{lastpage}
\usepackage{enumitem}
\usepackage{titlesec}
\usepackage{xcolor}
\usepackage{hyperref}
\usepackage{bookmark}
\usepackage[most]{tcolorbox}
\usepackage{setspace}
\usepackage{array}

\hypersetup{
  colorlinks=true,
  linkcolor=blue!50!black,
  urlcolor=blue!50!black
}

\setstretch{1.13}
\setlength{\parindent}{2em}
\setlength{\parskip}{0.25em}
\allowdisplaybreaks
\setlist[enumerate]{itemsep=0.45em, topsep=0.35em, leftmargin=2.4em}
\setlist[itemize]{itemsep=0.25em, topsep=0.25em, leftmargin=1.9em}

\definecolor{mainblue}{RGB}{36,83,140}
\definecolor{lightblue}{RGB}{241,246,252}
\definecolor{lightgray}{RGB}{246,246,246}
\definecolor{rulegray}{RGB}{110,118,128}

\titleformat{\section}
  {\Large\bfseries\color{mainblue}}
  {\thesection.}
  {0.45em}
  {}
  [\vspace{0.2em}\color{rulegray}\titlerule]
\titlespacing*{\section}{0pt}{1.1em}{0.7em}

\pagestyle{fancy}
\fancyhf{}
\fancyhead[L]{\small 函数图像对称周测 A 卷}
\fancyhead[C]{\small 代数专题：基础与标准变式}
\fancyhead[R]{\small 刘文博}
\fancyfoot[C]{\small 第 \thepage 页 / 共 \pageref{LastPage} 页}
\renewcommand{\headrulewidth}{0.55pt}
\renewcommand{\footrulewidth}{0.35pt}

\newtcolorbox{infobox}[1]{
  breakable,
  colback=lightblue,
  colframe=mainblue,
  boxrule=0.7pt,
  arc=1mm,
  title=\textbf{#1},
  fonttitle=\bfseries
}

\newenvironment{solution}{\par\noindent\textbf{参考解答：}}{\hfill$\square$\par}
\newcommand{\answerspace}{\vspace{2.3cm}}
\newcommand{\biganswerspace}{\vspace{3.1cm}}

\begin{document}

\thispagestyle{empty}
\begin{center}
{\fontsize{22}{28}\selectfont\bfseries 函数图像对称周测 A 卷\par}
\vspace{0.4cm}
{\large 适合初中阶段函数专题基础检测\par}
\end{center}

\begin{infobox}{考试说明}
\begin{itemize}
  \item 测试时间：60--70 分钟
  \item 满分：100 分
  \item A 卷侧重标准图像对称与反向判断，先分清“关于 x 轴改外面，关于 y 轴改里面”。
\end{itemize}
\end{infobox}

\section{试题}

\begin{enumerate}
  \item （10 分）把函数
  \[
  y=x^2+1
  \]
  关于 x 轴对称，求新方程。
  \answerspace

  \item （12 分）把函数
  \[
  y=2x-3
  \]
  关于 y 轴对称，求新方程。
  \answerspace

  \item （12 分）把函数
  \[
  y=x^3
  \]
  关于原点对称，求新方程。
  \answerspace

  \item （12 分）已知函数
  \[
  y=-x^2+4,
  \]
  说出它是由
  \[
  y=x^2-4
  \]
  怎样对称得到的。
  \answerspace

  \item （12 分）已知函数
  \[
  y=(x-4)^2+2,
  \]
  写出它关于 y 轴对称后的方程。
  \answerspace

  \item （14 分）判断函数
  \[
  y=x^4+3x^2+1
  \]
  的图像是否关于 y 轴对称，并说明理由。
  \biganswerspace

  \item （14 分）判断函数
  \[
  y=x^5-2x
  \]
  的图像是否关于原点对称，并说明理由。
  \biganswerspace

  \item （14 分）把函数
  \[
  y=|x-2|+3
  \]
  关于 y 轴对称，求新方程，并写出新图像顶点坐标。
  \biganswerspace
\end{enumerate}

\newpage
\section{详细解答}

\begin{enumerate}
  \item \begin{solution}
  关于 x 轴对称，就是在原函数外面加负号，所以
  \[
  \boxed{y=-(x^2+1)=-x^2-1}.
  \]
  \end{solution}

  \item \begin{solution}
  关于 y 轴对称，应把 $x$ 换成 $-x$，得到
  \[
  y=2(-x)-3=-2x-3.
  \]
  所以新方程为
  \[
  \boxed{y=-2x-3}.
  \]
  \end{solution}

  \item \begin{solution}
  关于原点对称，应写成
  \[
  y=-f(-x).
  \]
  这里 $f(x)=x^3$，所以
  \[
  y=-(-x)^3=x^3.
  \]
  因此新方程仍然是
  \[
  \boxed{y=x^3}.
  \]
  \end{solution}

  \item \begin{solution}
  原函数是
  \[
  y=x^2-4.
  \]
  关于 x 轴对称后，应变成
  \[
  y=-(x^2-4)=-x^2+4.
  \]
  因此题目中的函数是由原函数\textbf{关于 x 轴对称}得到的。
  \end{solution}

  \item \begin{solution}
  关于 y 轴对称时，应把式子里的 $x$ 换成 $-x$：
  \[
  y=((-x)-4)^2+2=(x+4)^2+2.
  \]
  所以新方程为
  \[
  \boxed{y=(x+4)^2+2}.
  \]
  \end{solution}

  \item \begin{solution}
  把 $x$ 换成 $-x$：
  \[
  f(-x)=(-x)^4+3(-x)^2+1=x^4+3x^2+1=f(x).
  \]
  因此它满足
  \[
  f(-x)=f(x),
  \]
  所以图像\textbf{关于 y 轴对称}。
  \end{solution}

  \item \begin{solution}
  把 $x$ 换成 $-x$：
  \[
  f(-x)=(-x)^5-2(-x)=-x^5+2x=-(x^5-2x)=-f(x).
  \]
  因此它满足
  \[
  f(-x)=-f(x),
  \]
  所以图像\textbf{关于原点对称}。
  \end{solution}

  \item \begin{solution}
  关于 y 轴对称，应把 $x$ 换成 $-x$：
  \[
  y=|(-x)-2|+3=|x+2|+3.
  \]
  所以新方程为
  \[
  \boxed{y=|x+2|+3}.
  \]

  原函数 $y=|x-2|+3$ 的顶点是 $(2,3)$，关于 y 轴对称后变成
  \[
  \boxed{(-2,3)}.
  \]
  \end{solution}
\end{enumerate}

\end{document}
